\documentclass[tikz,border=4pt]{standalone}
\input{common}
\usetikzlibrary{arrows.meta,positioning,calc}
\begin{document}
\begin{tikzpicture}[
    box/.style={draw,minimum height=7mm,inner sep=0pt},
    ptr/.style={-{Stealth[length=2.2mm]},thick},
    note/.style={font=\small,text=black!70}]
  % 结点 = 数据域 + 指针域
  \newcommand{\cell}[3]{%
    \node[box,minimum width=10mm] (#1d) at (#2,0) {#3};
    \node[box,minimum width=6mm,right=0pt of #1d] (#1p) {};}
  \node (head) at (-1.6,0) {\texttt{head}};
  \cell{a}{0}{$a_1$}
  \cell{b}{2.6}{$a_2$}
  \cell{c}{5.2}{$a_3$}
  \node at (7.1,0) {$\cdots$};
  \cell{z}{8.6}{$a_n$}
  \node at (zp) {$\wedge$};
  \foreach \n in {a,b,c} {\fill (\n p.center) circle (1.4pt);}
  \draw[ptr] (head) -- (ad.west);
  \draw[ptr] (ap.center) -- (bd.west);
  \draw[ptr] (bp.center) -- (cd.west);
  \draw[ptr] (cp.center) -- (6.6,0);
  \draw[ptr] (7.5,0) -- (zd.west);
  % 域的标注
  \draw[note,<-] (ad.north) -- ++(0,0.6) node[above] {数据域 \texttt{data}};
  \draw[note,<-] (ap.north) -- ++(0.9,0.6) node[above right=-2pt] {指针域 \texttt{next}};
  \node[note,below=2mm of zp] {$\nil$};
\end{tikzpicture}
\end{document}
